\chapter[El Poder Malgastado]{El Poder Malgastado\\ \large La autoridad sin compasión es un préstamo con intereses de miseria.}

\elegantimage{12_peso_injusticia/web/assets/hero.png}

\lettrine[lines=3, nindent=0.5em, findent=0.2em]{E}{}xplotar a los trabajadores para enriquecerse es una violación de la ley de equilibrio. El éxito depende de las manos de quienes desprecias.

\section{El Patrón de Arcilla}
\begin{elegantquote}
\textit{Era un jefe duro que robaba el sudor de sus jornaleros, pagándoles con humillación y migajas. Mientras él dormía en sábanas de seda fina, ellos apenas tenían cómo alimentar a sus familias. Se creía astuto y próspero...}

\textit{Pero la balanza invisible del universo mide la justicia con exactitud. De la noche a la mañana, sus riquezas se esfumaron. Se vio en la calle, mendigando un rincón oscuro donde descansar, sintiendo el desprecio y el rigor de quienes pasaban a su lado.}

\textit{Tuvo que ensuciar sus manos con el barro, sudar sangre por unas pocas monedas, y sentir la impotencia quemándole el pecho. Era necesario que el amo se hiciera esclavo, para que entendiera que el sustento es sagrado, que el trabajador merece honra... y que la verdadera riqueza es prosperar sin pisar jamás a nadie.}

\end{elegantquote}

\vspace{1em}
\begin{center}
    \textbf{\Large Nadie es tan rico que no pueda necesitar ayuda.}
\end{center}
\vspace{1em}

\section{Análisis Kármico y Traducción}
\begin{wrapfigure}{r}{0.5\textwidth}
  \vspace{-1.5em}
  \begin{center}
  \inlineelegantimage[0.45\textwidth]{12_peso_injusticia/web/assets/pasaje_original.jpg}
  \end{center}
  \vspace{-1.5em}
\end{wrapfigure}
Tratar con injusticia y maldad a quienes nos sirven consume velozmente nuestras bendiciones materiales, hundiendo el destino en la ruina para enseñarnos el respeto sagrado hacia el trabajador.

